\chapter{Метод гібридного постквантового наскрізного шифрування}
\label{ch:e2ee}

% ============================================================
%  РОЗДІЛ 6 -- криптографічний контур, частина 2: обмін повідомленнями.
%  Джерело: aura-protected-protocol-rs.
%  УВАГА: закрити documentation drift по ProVerif (3/6 vs 4/4 vs 6/6).
% ============================================================

\section{Постановка та гібридний криптографічний корінь}
\label{sec:e2ee-root}
% -- X25519 + ML-KEM-768 як гібридний корінь [cite:fips203];
%    мотивація гібрида в перехідній моделі.

\section{Класична автентифікація сторін}
\label{sec:e2ee-auth}
% -- Ed25519; ЯВНО зазначити: це класичний підпис, НЕ постквантовий.

\section{Подвійний «храповик» з KEM-підсиленим кроком}
\label{sec:e2ee-ratchet}
% -- Double Ratchet [cite:signal-doubleratchet] + KEM-binding на рівні ratchet.

\section{Захист метаданих}
\label{sec:e2ee-metadata}
% -- AEAD/MRAE payload protection; ротація ключів метаданих.

\section{Межі: replay, rollback, skipped-key}
\label{sec:e2ee-limits}
% -- Обмеження та контрзаходи; межі захисту.

\section{Формальні докази та узгодження тверджень}
\label{sec:e2ee-formal}
% -- Tamarin [cite:tamarin] / ProVerif [cite:proverif], scoped.
% -- ЧЕСНЕ формулювання (узгоджено з межами тверджень, розділ 2):
%    ProVerif розряджає 3 первинні твердження + 1 додаткову scoped-перевірку
%    чесного повідомлення; груповий scoped-модель -- 6 запитів епохи групи;
%    stress-зобов'язання залишаються non-claims.
Формальні гарантії подано у scoped-формі: активна модель ProVerif розряджає три первинні твердження та одну додаткову перевірку кореспонденції чесного повідомлення, а stress-зобов'язання явно не належать до тверджень. \emph{[Звірити з логами та статтею перед фіналом.]}

\section{Груповий протокол (scoped)}
\label{sec:e2ee-group}
% -- Як підрозділ; груповий scoped-модель 6/6; НЕ називати повним async MLS-доказом.

\rozdilconcl{6}
\emph{[Висновки до розділу 6.]}
